%%%%%%%%%%%%%%%%%%%%%%%%%%%%%%%%%%%%%%%%%%%%%%%%%%%%%%%%%%%%%%%%%%%%%%
% Lớp 8 | Chương I. Đa thức | Phép nhân đa thức
% Dạng 8D1?4-1: Nhân đơn thức với đa thức
%%%%%%%%%%%%%%%%%%%%%%%%%%%%%%%%%%%%%%%%%%%%%%%%%%%%%%%%%%%%%%%%%%%%%%


% nguồn: De_AI_Toan8_cuoi_hoc_ki_i_1 (nhập 21/11/2026)
\begin{bt}%[8D1H4-1]%[ma:9433dbfb]
	Rút gọn biểu thức $A = 2x(x^2 - 3x + 1) - x^2(2x - 6)$ rồi tính giá trị của $A$ tại $x = -15$.
	\loigiai{
		Ta có $A = (2x^3 - 6x^2 + 2x) - (2x^3 - 6x^2) = 2x^3 - 6x^2 + 2x - 2x^3 + 6x^2 = 2x$. \\
Thay $x = -15$ vào $A$, ta được $A = 2 \cdot (-15) = -30$.
	}
\end{bt}

% nguồn: De_AI_Toan8_cuoi_hoc_ki_i_1 (nhập 21/11/2026)
% nguồn: De_AI_Toan8_cuoi_hoc_ki_i_1 (nhập 21/11/2026)
\begin{bt}%[8D1H4-1]%[ma:4dbcb615]
	Chứng minh rằng giá trị của biểu thức $C = 3x(x^2 - 2y) - x^2(3x - y) + xy(6 - x)$ luôn bằng $0$ với mọi giá trị của $x$ và $y$.
	\loigiai{
		Ta có $C = (3x^3 - 6xy) - (3x^3 - x^2y) + (6xy - x^2y)$ \\
$= 3x^3 - 6xy - 3x^3 + x^2y + 6xy - x^2y = 0$. \\
Vậy giá trị của $C$ luôn bằng $0$ với mọi $x, y$.
	}
\end{bt}

% nguồn: chuyen_dang (nhập 22/11/2026)
\begin{bt}%[8D1H4-1]%[ma:b95fa86c]
	Rút gọn biểu thức $A = (x + 2)(x^2 - 2x + 4) - (x - 1)(x^2 + x + 1)$ rồi chứng tỏ giá trị của $A$ không phụ thuộc vào $x$.
	\loigiai{
		Áp dụng hằng đẳng thức tổng hai lập phương và hiệu hai lập phương: \\
$(x + 2)(x^2 - 2x + 4) = x^3 + 2^3 = x^3 + 8$. \\
$(x - 1)(x^2 + x + 1) = x^3 - 1^3 = x^3 - 1$. \\
Suy ra $A = (x^3 + 8) - (x^3 - 1) = x^3 + 8 - x^3 + 1 = 9$. \\
Vậy giá trị của $A$ luôn bằng $9$, không phụ thuộc vào biến $x$.
	}
\end{bt}
